\begin{lemma}[Case 2, observation (21)]
Put \(X=V(T_n)\cup V(T'_j)\), and suppose
\(v_n\in V(T_n)\cap V(T'_j)\).  Under \(\varrho_1\), let \(P\) be a
simple \((\gamma_n,\delta_n)\)-path in the incidence graph from \(v_n\)
to \(x\) whose support is the entire
\(\noderef{TwoColorComponent}\) through \(v_n\).  Assume \(x\notin X\)
and that \(P\) contains a vertex of \(T_n\) other than \(v_n\).  Suppose
\[
 \partial_{\varrho_1,\gamma_n}(T_n)=\{f_n\},\qquad
 \partial_{\varrho_1,\delta_n}(X)=\varnothing,                 \tag{20}
\]
and \(f_n\) is incident with \(v_n\).  Then there are a vertex \(u\) and
a simple subpath \(Q=P[x,u]\), represented as an exact incidence-graph
walk from \(x\) to \(u\), such that
\[
 u\in V(T'_j)\setminus V(T_n),                                 \tag{21}
\]
\(Q\) uses at least one graph-edge identity, every node of \(Q\) is a
node of \(P\), and \(u\) is the only vertex-node of \(Q\) in \(X\).
Moreover, \(Q\) contains an exact \(\gamma_n\)-colored boundary identity
of \(X\) incident with \(u\), and \(u\) is not incident with \(f_n\).
\end{lemma}
\begin{proof}
Write the nodes of the simple walk \(P\) in their exact walk order.  The
first vertex-node \(v_n\) lies in \(X\), whereas the last vertex-node
\(x\) does not.  Among the indices of vertex-nodes of \(P\) whose
vertices lie in \(X\), choose the largest one, and call its vertex \(u\).
The finite maximum exists because the first index is eligible.  Reverse
the suffix from that occurrence of \(u\) to \(x\); this is the displayed
walk \(Q=P[x,u]\).  A contiguous subwalk of a simple walk is simple, and
its nodes are nodes of \(P\).  Maximality of the chosen occurrence says
exactly that \(u\) is the only vertex-node of \(Q\) lying in \(X\).
Since \(x\notin X\), the suffix is nonconstant; the bipartite adjacency
rule of \noderef{TwoColorIncidenceGraph} makes its incidence-graph length
at least two.

Let \(g\) be the first graph-edge node of the unreversed suffix, namely
the exact edge identity by which \(P\) leaves \(X\) at \(u\).  Its next
vertex-node is outside \(X\), so \noderef{BoundaryEdges} puts
\(g\in\partial_G(X)\), and \noderef{Incident} says that \(g\) is
incident with \(u\).  Adjacency in
\noderef{TwoColorIncidenceGraph} says that its
\(\varrho_1\)-color is either \(\gamma_n\) or \(\delta_n\).  The second
equation in (20) excludes \(\delta_n\); hence
\[
             g\in\partial_{\varrho_1,\gamma_n}(X).             \tag{1}
\]
We show that \(u\notin V(T_n)\).  First, the occurrence of \(u\) is not
the initial occurrence of \(v_n\): by hypothesis there is a distinct
vertex \(w\in V(T_n)\subseteq X\) later on \(P\), and the largest
in-\(X\) occurrence is at least as late as that of \(w\).  Thus
\(u\ne v_n\).

Suppose nevertheless that \(u\in V(T_n)\).  The endpoint after \(g\) is
outside \(X\), and hence outside \(V(T_n)\), so \(g\) is also a boundary
identity of \(T_n\).  Equation (1) and the first equality in (20) give
the exact identity \(g=f_n\).  Both \(u\) and \(v_n\) are incident with
this loopless identity.  Since the other endpoint of \(g\) is outside
\(X\) while \(u,v_n\in X\), the two-endpoint clause of
\noderef{LooplessMultigraph} forces \(u=v_n\), contrary to the preceding
paragraph.  Therefore \(u\notin V(T_n)\).  Its defining membership in
\(X=V(T_n)\cup V(T'_j)\) now gives \(u\in V(T'_j)\), proving (21).

Finally suppose that \(u\) were incident with \(f_n\).  The singleton
identity in (20) says that \(f_n\) has color \(\gamma_n\), and its
incidence with \(v_n\) makes its edge-node adjacent to the initial node
of \(P\) in \noderef{TwoColorIncidenceGraph}.  It therefore belongs to
the component through \(v_n\), and the full-support hypothesis puts that
edge-node on \(P\).  An edge-node has at most its two actual endpoint
neighbors by \noderef{LooplessMultigraph}.  Since \(P\) is a simple path
whose endpoints are vertex-nodes, both neighbors used by \(P\) at the
interior edge-node \(f_n\) are its two endpoints, namely \(v_n\) and
\(u\).  Thus \(u\) is the first vertex of \(P\) after \(v_n\).  The
distinct \(T_n\)-vertex supplied by the two-vertex hypothesis occurs
later, so the last vertex of \(P\) in \(X\) cannot be \(u\), contrary to
the definition of \(u\).  Thus \(u\) is not incident with \(f_n\), and
all asserted exact path and edge data have been retained.
\end{proof}
